\subsection{Foncteurs hypertor locaux de complexes de Modules quasi-cohérents : cas des schémas affines.}

\begin{env}[6.4.1]
Soient $S$ un schéma affine d'anneau $A$, $X$, $Y$ deux $S$-schémas affines
d'anneaux $B$, $C$ respectivement, de sorte que $B$ et $C$ sont des algèbres
sur $A$. Tout complexe $\mathscr P_\bullet$ (resp. $\mathscr Q_\bullet$) de
$\mathscr O_X$-Modules (resp. $\mathscr O_Y$-Modules) quasi-cohérents est de
la forme $\widetilde{P}_\bullet$ (resp. $\widetilde{Q}_\bullet$), où
$P_\bullet$ (resp. $Q_\bullet$) est un complexe de $B$-modules (resp.
$C$-modules) (\textbf{I}, 1.3.7 et 1.3.8). On peut évidemment considérer $P_\bullet$
et $Q_\bullet$ comme des complexes de $A$-modules et former les
$\mathbf{Tor}_n^A(P_\bullet,Q_\bullet)$; en outre, en vertu du caractère
bifonctoriel de $\mathbf{Tor}_n^A(P_\bullet,Q_\bullet)$, les $A$-algèbres
$B$ et $C$ opèrent dans ce $A$-module, et ces opérations en font un
$(B,C)$-bimodule, ou, ce qui revient au même, un module sur
$B\otimes_A C=A(X\times_S Y)$. On a donc défini de la sorte un
$\mathscr O_{X\times_S Y}$-Module quasi-cohérent
\[
  \boldsymbol{\mathscr{T}\!or}_n^{\mathscr O_S}(\mathscr P_\bullet,\mathscr Q_\bullet)
  =\bigl(\mathbf{Tor}_n^A(P_\bullet,Q_\bullet)\bigr)^{\sim}
  \tag{6.4.1.1}\label{III.6.4.1.1.fr}
\]
que l'on appelle l'\emph{hypertor local} d'indice $n$ des complexes
$\mathscr P_\bullet$ et $\mathscr Q_\bullet$ et que l'on note aussi
$\boldsymbol{\mathscr{T}\!or}_n^S(\mathscr P_\bullet,\mathscr Q_\bullet)$.
\end{env}

\begin{namedtheorem}[6.4.2]{Lemme}
Avec les notations de \textup{(6.4.1)}, supposons que l'anneau $A$ soit de la forme
$R^{-1}A'$, où $A'$ est un anneau et $R$ une partie multiplicative de $A'$.
Soit $S'=\operatorname{Spec}(A')$, de sorte que $X$ et $Y$ peuvent être
considérés comme des $S'$-préschémas et que l'on a
$X\times_{S'}Y=X\times_S Y$ \textup{(\textbf{I}, 1.6.2} et \textup{3.2.4)}. On a alors
$\boldsymbol{\mathscr{T}\!or}_\bullet^{S'}(\mathscr P_\bullet,\mathscr Q_\bullet)
=\boldsymbol{\mathscr{T}\!or}_\bullet^S(\mathscr P_\bullet,\mathscr Q_\bullet)$.
\end{namedtheorem}

\begin{proof}
Cela résulte de la formule (6.4.1.1) et de (6.3.13).
\end{proof}

\begin{env}[6.4.3]
Avec les notations et hypothèses de (6.4.1), soient
$\mathscr F=\widetilde F$ un $\mathscr O_X$-Module quasi-cohérent,
$\mathscr G=\widetilde G$ un $\mathscr O_Y$-Module quasi-cohérent;
considérant $\mathscr F$ et $\mathscr G$ comme des complexes de Modules, on
notera $\mathscr{T}\!or_n^{\mathscr O_S}(\mathscr F,\mathscr G)$ ou
$\mathscr{T}\!or_n^S(\mathscr F,\mathscr G)$ leur hypertor d'indice $n$; il
résulte de (6.3.5 (ii)) que l'on a
\[
  \mathscr{T}\!or_n^{\mathscr O_S}(\mathscr F,\mathscr G)
  =\bigl(\operatorname{Tor}_n^A(F,G)\bigr)^{\sim}.
  \tag{6.4.3.1}\label{III.6.4.3.1.fr}
\]

Revenons alors au cas général de deux complexes de Modules quasi-cohérents
$\mathscr P_\bullet$, $\mathscr Q_\bullet$. Les formules (6.4.1.1) et
(6.4.3.1) montrent, compte tenu de la prop. (6.3.2), que
$\boldsymbol{\mathscr{T}\!or}_\bullet^S(\mathscr P_\bullet,\mathscr Q_\bullet)$ est
l'aboutissement de deux suites spectrales
$'\mathscr E(\mathscr P_\bullet,\mathscr Q_\bullet)$,
$''\mathscr E(\mathscr P_\bullet,\mathscr Q_\bullet)$, dont les termes
$E_2$ sont donnés par
\[
  '\mathscr E_{pq}^2
  =\mathscr H_p\bigl(\mathscr{T}\!or_q^S
    (\mathscr P_\bullet,\mathscr Q_\bullet)\bigr)
  \tag{6.4.3.2}\label{III.6.4.3.2.fr}
\]
\[
  ''\mathscr E_{pq}^2
  =\bigoplus_{q'+q''=q}
  \mathscr{T}\!or_p^S
  \bigl(\mathscr H_{q'}(\mathscr P_\bullet),
        \mathscr H_{q''}(\mathscr Q_\bullet)\bigr)
  \tag{6.4.3.3}\label{III.6.4.3.3.fr}
\]
où $\mathscr{T}\!or_q^S(\mathscr P_\bullet,\mathscr Q_\bullet)$ est le
bicomplexe de $\mathscr O_{X\times_S Y}$-Modules quasi-cohérents
$\mathscr{T}\!or_q^S(\mathscr P_i,\mathscr Q_j)$.
\end{env}

\begin{env}[6.4.4]
Considérons maintenant deux autres schémas affines
$X^{(1)}=\operatorname{Spec}(B^{(1)})$,
$Y^{(1)}=\operatorname{Spec}(C^{(1)})$, où $B^{(1)}$ et $C^{(1)}$ sont des
$A$-algèbres, et supposons donnés deux
\oldpage[III]{147}
S-morphismes $u:X^{(1)}\to X$, $v:Y^{(1)}\to Y$, correspondant à des
$A$-homomorphismes $\varphi:B\to B^{(1)}$, $\psi:C\to C^{(1)}$.
Considérons les complexes
$u^*(\mathscr P_\bullet)=(\mathscr P_\bullet\otimes_B B^{(1)})^{\sim}$ de
$\mathscr O_{X^{(1)}}$-Modules,
$v^*(\mathscr Q_\bullet)=(Q_\bullet\otimes_C C^{(1)})^{\sim}$ de
$\mathscr O_{Y^{(1)}}$-Modules. Les $A$-homomorphismes canoniques
\[
  P_\bullet\longrightarrow P_\bullet\otimes_B B^{(1)},
  \qquad
  Q_\bullet\longrightarrow Q_\bullet\otimes_C C^{(1)}
\]
donnent par fonctorialité un $A$-homomorphisme
\[
  \mathbf{Tor}_\bullet^A(P_\bullet,Q_\bullet)
  \longrightarrow
  \mathbf{Tor}_\bullet^A
  (P_\bullet\otimes_B B^{(1)},Q_\bullet\otimes_C C^{(1)});
\]
en outre, toujours par fonctorialité, cet homomorphisme est en fait un
homomorphisme de $(B\otimes_A C)$-modules. D'où l'on conclut que l'on a défini
ainsi un $(u\times_S v)$-morphisme
\[
  \theta:\boldsymbol{\mathscr{T}\!or}_\bullet^S(\mathscr P_\bullet,\mathscr Q_\bullet)
  \longrightarrow
  \boldsymbol{\mathscr{T}\!or}_\bullet^S
  (u^*(\mathscr P_\bullet),v^*(\mathscr Q_\bullet))
  \tag{6.4.4.1}\label{III.6.4.4.1.fr}
\]
et par suite, un homomorphisme de
$\mathscr O_{X^{(1)}\times_S Y^{(1)}}$-Modules
\[
  \theta^\sharp:(u\times_S v)^*
  \bigl(\boldsymbol{\mathscr{T}\!or}_\bullet^S(\mathscr P_\bullet,\mathscr Q_\bullet)\bigr)
  \longrightarrow
  \boldsymbol{\mathscr{T}\!or}_\bullet^S
  (u^*(\mathscr P_\bullet),v^*(\mathscr Q_\bullet))
  \tag{6.4.4.2}\label{III.6.4.4.2.fr}
\]
qui est évidemment un morphisme de \emph{bi-$\partial$-foncteurs} dans les catégories
de Modules quasi-cohérents limités inférieurement.

L'homomorphisme (6.4.4.2) n'est pas nécessairement bijectif; toutefois :
\end{env}

\begin{namedtheorem}[6.4.5]{Lemme}
Avec les notations de \textup{(6.4.4)}, supposons que $u$ et $v$ soient des immersions
ouvertes; alors l'homomorphisme \textup{(6.4.4.2)} est bijectif.
\end{namedtheorem}

\begin{proof}
Identifions $X^{(1)}$ (resp. $Y^{(1)}$) à un ouvert de $X$ (resp. $Y$);
$X^{(1)}$ (resp. $Y^{(1)}$) est alors réunion d'ouverts de la forme
$\mathrm D(f)$ (resp. $\mathrm D(g)$), où $f\in B$ (resp. $g\in C$), et les
préschémas induits $\mathrm D(\varphi(f))$ et $\mathrm D(f)$ (resp.
$\mathrm D(\psi(g))$ et $\mathrm D(g)$) sont isomorphes. Il suffira de prouver
le lemme lorsque $X^{(1)}$ (resp. $Y^{(1)}$) est de la forme $\mathrm D(f)$
(resp. $\mathrm D(g)$); en effet, si ce point est établi, et si l'on revient au
cas général, il suffira de prouver que la restriction de $\theta^\sharp$ à
chaque ouvert $\mathrm D(f)\times_S\mathrm D(g)$ est un isomorphisme; or, si
$u_1:\mathrm D(f)\to X^{(1)}$, $v_1:\mathrm D(g)\to Y^{(1)}$ sont les
injections canoniques, la restriction précédente n'est autre que
$(u_1\times_S v_1)^*(\theta^\sharp)$; mais il est immédiat, en vertu des
définitions (6.4.4) et de ($\mathbf{0}_{\mathrm I}$, 4.4.8), qu'en la composant avec
l'homomorphisme canonique
\[
  (u_1\times_S v_1)^*
  \bigl(\boldsymbol{\mathscr{T}\!or}_\bullet^S
    (u^*(\mathscr P_\bullet),v^*(\mathscr Q_\bullet))\bigr)
  \longrightarrow
  \boldsymbol{\mathscr{T}\!or}_\bullet^S
  (u'^*(\mathscr P_\bullet),v'^*(\mathscr Q_\bullet))
  \tag{6.4.5.1}\label{III.6.4.5.1.fr}
\]
où $u'=u\circ u_1$ et $v'=v\circ v_1$, on obtient l'homomorphisme canonique
\[
  (u'\times_S v')^*
  \bigl(\boldsymbol{\mathscr{T}\!or}_\bullet^S(\mathscr P_\bullet,\mathscr Q_\bullet)\bigr)
  \longrightarrow
  \boldsymbol{\mathscr{T}\!or}_\bullet^S
  (u'^*(\mathscr P_\bullet),v'^*(\mathscr Q_\bullet))
  \tag{6.4.5.2}\label{III.6.4.5.2.fr}
\]
et si l'on sait que (6.4.5.1) et (6.4.5.2) sont des isomorphismes, il en
résultera qu'il en est de même de $(u_1\times_S v_1)^*(\theta^\sharp)$.

Supposons donc que $X^{(1)}=\mathrm D(f)$ et $Y^{(1)}=\mathrm D(g)$, de sorte
que $B^{(1)}=B_f$ et $C^{(1)}=C_g$; $u^*(\mathscr P_\bullet)$ (resp.
$v^*(\mathscr Q_\bullet)$) s'identifie alors à
$(P_\bullet)_f^{\sim}$ (resp. $(Q_\bullet)_g^{\sim}$); d'autre part,
$X^{(1)}\times_S Y^{(1)}$ s'identifie au sous-schéma ouvert
$\mathrm D(f\otimes g)$ de
$X\times_S Y=\operatorname{Spec}(B\otimes_A C)$ (\textbf{II}, 4.3.2.4); il s'agit de
prouver que l'homomorphisme
\[
  \bigl(\mathbf{Tor}_\bullet^A(P_\bullet,Q_\bullet)\bigr)_{f\otimes g}
  \longrightarrow
  \mathbf{Tor}_\bullet^A((P_\bullet)_f,(Q_\bullet)_g)
  \tag{6.4.5.3}\label{III.6.4.5.3.fr}
\]
\oldpage[III]{148}
déduit par fonctorialité des homomorphismes canoniques
$P_\bullet\to(P_\bullet)_f$, $Q_\bullet\to(Q_\bullet)_g$, est bijectif. Or
($\mathbf{0}_{\mathrm I}$, 1.6.1), on peut écrire
$(P_\bullet)_f=\varinjlim P_\bullet^{(n)}$, où les $P_\bullet^{(n)}$ sont tous
des complexes de $B$-modules identiques à $P_\bullet$, l'application
$P_\bullet^{(m)}\to P_\bullet^{(n)}$ pour $m\leqslant n$ étant la
multiplication par $f^{n-m}$; on a un résultat analogue pour $Q_\bullet$ en
remplaçant $f$ par $g$; d'autre part, il est clair que l'homomorphisme
\[
  \mathbf{Tor}_\bullet^A(P_\bullet^{(m)},Q_\bullet^{(m)})
  \longrightarrow
  \mathbf{Tor}_\bullet^A(P_\bullet^{(n)},Q_\bullet^{(n)})
\]
correspondant aux homomorphismes
$P_\bullet^{(m)}\to P_\bullet^{(n)}$ et
$Q_\bullet^{(m)}\to Q_\bullet^{(n)}$ est par définition la multiplication
par $(f\otimes g)^{n-m}$. La conclusion résulte alors de ($\mathbf{0}_{\mathrm I}$,
1.6.1) appliqué au premier membre de (6.4.5.3) et de (6.3.6).
\end{proof}

\begin{env}[6.4.6]
Avec les notations de (6.4.4), on définit de même des homomorphismes canoniques
de foncteurs spectraux
\[
  \left\{
  \begin{aligned}
  (u\times_S \boldsymbol{v})^*\bigl('\mathscr E(\mathscr P_\bullet,\mathscr Q_\bullet)\bigr)
    &\longrightarrow
      '\mathscr E(u^*(\mathscr P_\bullet),v^*(\mathscr Q_\bullet))\\
  (u\times_S v)^*\bigl(''\mathscr E(\mathscr P_\bullet,\mathscr Q_\bullet)\bigr)
    &\longrightarrow
      ''\mathscr E(u^*(\mathscr P_\bullet),v^*(\mathscr Q_\bullet))
  \end{aligned}
  \right.
  \tag{6.4.6.1}\label{III.6.4.6.1.fr}
\]
et le raisonnement de (6.4.5) montre que lorsque $u$ et $v$ sont des
\emph{immersions ouvertes}, les homomorphismes (6.4.6.1) sont \emph{bijectifs} : en effet,
compte tenu de (6.3.6.2) et (6.3.6.3), il prouve que c'est un isomorphisme pour
les termes $E^2$, et (6.4.5) montre que c'est un isomorphisme pour les
aboutissements; on conclut donc à l'aide de (\textbf{0}, 11.1.2) et (\textbf{0}, 11.2.4).
\end{env}

\subsection{Foncteurs hypertor locaux de complexes de Modules quasi-cohérents : cas général.}

\begin{env}[6.5.1]
Considérons maintenant un préschéma $S$ quelconque et deux $S$-préschémas
quelconques $X$, $Y$; soit $\mathscr P_\bullet$ (resp. $\mathscr Q_\bullet$)
un complexe de $\mathscr O_X$-Modules (resp. $\mathscr O_Y$-Modules)
quasi-cohérents. Posons $Z=X\times_S Y$; nous allons définir des
$\mathscr O_Z$-Modules quasi-cohérents
$\boldsymbol{\mathscr{T}\!or}_n^S(\mathscr P_\bullet,\mathscr Q_\bullet)$ dits
\emph{hypertor locaux} de $\mathscr P_\bullet$ et $\mathscr Q_\bullet$ qui se
réduiront à ceux déjà définis dans (6.4) lorsque $S$, $X$ et $Y$ sont affines.

Lorsque $\mathscr P_\bullet$ et $\mathscr Q_\bullet$ se réduisent
respectivement à leurs termes de degré 0, $\mathscr F$ et $\mathscr G$, (les
autres étant nuls), on écrira
$\mathscr{T}\!or_n^S(\mathscr F,\mathscr G)$ au lieu de
$\boldsymbol{\mathscr{T}\!or}_n^S(\mathscr P_\bullet,\mathscr Q_\bullet)$.
\end{env}

\begin{env}[6.5.2]
Supposons d'abord $S$ \emph{affine}, et soient $(X_\lambda)$, $(Y_\mu)$ des
recouvrements de $X$ et $Y$ respectivement par des ouverts affines; alors les
$Z_{\lambda\mu}=X_\lambda\times_S Y_\mu$ forment un recouvrement ouvert affine
de $Z$. Posons
$\mathscr P_{\lambda\bullet}=\mathscr P_\bullet|X_\lambda$,
$\mathscr Q_{\mu\bullet}=\mathscr Q_\bullet|Y_\mu$; nous avons donc pour tout
couple $(\lambda,\mu)$ un $\mathscr O_{Z_{\lambda\mu}}$-Module quasi-cohérent
$\mathscr F_{\lambda\mu}
  =\boldsymbol{\mathscr{T}\!or}_n^S
   (\mathscr P_{\lambda\bullet},\mathscr Q_{\mu\bullet})$,
et il faut montrer que les $\mathscr F_{\lambda\mu}$ vérifient la condition de
recollement ($\mathbf{0}_{\mathrm I}$, 3.3.1). Pour cela, il suffit de vérifier que pour
tout ouvert affine $U\subset X_\lambda\cap X_{\lambda'}$ (resp.
$V\subset Y_\mu\cap Y_{\mu'}$), les restrictions de
$\mathscr F_{\lambda\mu}$ et de $\mathscr F_{\lambda'\mu'}$ à $U\times_S V$
sont canoniquement isomorphes; mais cela découle aussitôt de l'existence
d'isomorphismes canoniques de ces restrictions sur
$\boldsymbol{\mathscr{T}\!or}_n^S(\mathscr P_\bullet|U,\mathscr Q_\bullet|V)$ (6.4.5). En
outre, il résulte aussitôt de cette définition et de (6.4.5) que le
$\mathscr O_Z$-\textbf{Module} ainsi défini ne dépend pas (à un isomorphisme près) des
recouvrements ouverts $\boldsymbol{(}X_\lambda)$, $(Y_\mu)$ consi\oldpage[III]{149}dérés;
nous le noterons donc
$\boldsymbol{\mathscr{T}\!or}_n^S(\mathscr P_\bullet,\mathscr Q_\bullet)$; il résulte enfin
de (6.4.5) que pour \emph{toute} partie ouverte $U$ (resp. $V$) de $X$ (resp. $Y$), la
restriction de $\boldsymbol{\mathscr{T}\!or}_n^S(\mathscr P_\bullet,\mathscr Q_\bullet)$ à
$U\times_S V$ est canoniquement isomorphe à
$\boldsymbol{\mathscr{T}\!or}_n^S(\mathscr P_\bullet|U,\mathscr Q_\bullet|V)$.
\end{env}

\begin{env}[6.5.3]
Passons maintenant au cas général où $S$ est quelconque, \textbf{et} soit $(S_\alpha)$
un recouvrement de $S$ formé d'ouverts affines; désignons par $X_\alpha$
(resp. $Y_\alpha$) l'image réciproque \textbf{de} $S_\alpha$ dans $X$ (resp. $Y$); il
faut encore prouver que les faisceaux
$\boldsymbol{\mathscr{T}\!or}_n^{S_\alpha}
  (\mathscr P_\bullet|X_\alpha,\mathscr Q_\bullet|Y_\alpha)
  =\mathscr G_\alpha$
vérifient la condition de recollement. Il suffit de définir, pour tout ouvert
affine $T$ contenu dans $S_\alpha\cap S_\beta$, des isomorphismes canoniques des
restrictions de $\mathscr G_\alpha$ et $\mathscr G_\beta$ à $U\times_S V$ (en
désignant par $U$ et $V$ les images réciproques de $T$ dans $X$ et $Y$
respectivement) sur
$\boldsymbol{\mathscr{T}\!or}_n^T(\mathscr P_\bullet|U,\mathscr Q_\bullet|V)$; on peut, en
outre, se borner au cas où $T$ s'écrit à la fois $\mathrm D(f_\alpha)$ et
$\mathrm D(f_\beta)$, $f_\alpha$ (resp. $f_\beta$) étant une section de
$\mathscr O_S$ au-dessus de $S_\alpha$ (resp. $S_\beta$); mais alors
$\boldsymbol{\mathscr{T}\!or}_n^T(\mathscr P_\bullet|U,\mathscr Q_\bullet|V)$ est
canoniquement isomorphe à
$\boldsymbol{\mathscr{T}\!or}_n^{S_\alpha}(\mathscr P_\bullet|U,\mathscr Q_\bullet|V)$
d'une part, à
$\boldsymbol{\mathscr{T}\!or}_n^{S_\beta}(\mathscr P_\bullet|U,\mathscr Q_\bullet|V)$
d'autre part, en vertu de (6.4.2); comme on vient de définir des isomorphismes
canoniques de $\mathscr G_\alpha$ sur
$\boldsymbol{\mathscr{T}\!or}_n^{S_\alpha}(\mathscr P_\bullet|U,\mathscr Q_\bullet|V)$ et
de $\mathscr G_\beta$ sur
$\boldsymbol{\mathscr{T}\!or}_n^{S_\beta}(\mathscr P_\bullet|U,\mathscr Q_\bullet|V)$
(6.5.2), cela achève de définir le $\mathscr O_Z$-Module
$\boldsymbol{\mathscr{T}\!or}_n^S(\mathscr P_\bullet,\mathscr Q_\bullet)$. En outre, pour
toute partie ouverte $U$ (resp. $V$) de $X$ (resp. $Y$),
$\boldsymbol{\mathscr{T}\!or}_n^S(\mathscr P_\bullet|U,\mathscr Q_\bullet|V)$ est
canoniquement isomorphe à la restriction de
$\boldsymbol{\mathscr{T}\!or}_n^S(\mathscr P_\bullet,\mathscr Q_\bullet)$ à $U\times_S V$.

Il est immédiat que l'on a ainsi défini (dans les catégories de complexes de
Modules quasi-cohérents limités inférieurement) un \emph{bi-$\partial$-foncteur}
$\boldsymbol{\mathscr{T}\!or}_\bullet^S(\mathscr P_\bullet,\mathscr Q_\bullet)$ à valeurs
dans la catégorie des $\mathscr O_Z$-Modules, car il est clair que la question
est locale sur $X$, $Y$ et $S$, en vertu de (6.4.5) et de la remarque que
(6.4.4.2) est un morphisme de bi-$\partial$-foncteurs. On notera que si
$\mathscr P_\bullet$ et $\mathscr Q_\bullet$ sont réduits respectivement à
leurs termes de degré 0, $\mathscr F$ et $\mathscr G$,
$\mathscr{T}\!or_0^S(\mathscr F,\mathscr G)$ n'est autre, en vertu de
(6.4.1.1), que le produit tensoriel externe $\mathscr F\otimes_S\mathscr G$
défini dans (\textbf{I}, 9.1.2); cela résulte en effet de (\textbf{I}, 9.1.3).
\end{env}

\begin{env}[6.5.4]
Il résulte de la construction précédente et des remarques faites dans (6.4.6)
que $\boldsymbol{\mathscr{T}\!or}_\bullet^S(\mathscr P_\bullet,\mathscr Q_\bullet)$ est
l'aboutissement de deux foncteurs spectraux,
$'\mathscr E(\mathscr P_\bullet,\mathscr Q_\bullet)$,
$''\mathscr E(\mathscr P_\bullet,\mathscr Q_\bullet)$, de termes $E_2$ égaux à
\[
  '\mathscr E_{pq}^2
  =\mathscr H_p\bigl(\mathscr{T}\!or_q^S
    (\mathscr P_\bullet,\mathscr Q_\bullet)
  \tag{6.5.4.1}\label{III.6.5.4.1.fr}
\]
\[
  ''\mathscr E_{pq}^2
  =\bigoplus_{q'+q''=q}
   \mathscr{T}\!or_p^S
   \bigl(\mathscr H_{q'}(\mathscr P_\bullet),
         \mathscr H_{q''}(\mathscr Q_\bullet)\bigr)
  \tag{6.5.4.2}\label{III.6.5.4.2.fr}
\]

La suite spectrale (6.5.4.2) est toujours régulière; les deux suites spectrales
sont birégulières si $\mathscr P_\bullet$ et $\mathscr Q_\bullet$ sont \emph{limités
inférieurement}. Un autre cas où les deux suites précédentes sont birégulières
est le suivant :
\end{env}

\begin{env}[6.5.5]
Nous dirons que sur un espace topologique $T$ un faisceau d'anneaux
$\mathscr A$ est \emph{de dimension cohomologique $\leqslant n$} si, pour tout
$t\in T$ l'anneau $\mathscr A_t$ est de dimension cohomologique $\leqslant n$;
on dira alors aussi que l'\emph{espace annelé} $(T,\mathscr A)$ est
\emph{de dimension cohomologique $\leqslant n$}. On dira qu'un faisceau
d'anneaux (resp. un espace annelé) est de dimension cohomologique \emph{finie}
s'il existe un entier $n$ tel qu'il soit de dimension cohomologique
$\leqslant n$. On notera que si les $\mathscr A_t$ sont des \emph{anneaux locaux
(commutatifs) noethériens},
\oldpage[III]{150}
dire qu'ils sont de dimension cohomologique $\leqslant n$ signifie qu'ils sont
\emph{réguliers} et de \emph{dimension} (de Krull) $\leqslant n$ ($\mathbf{0}_{\mathrm{IV}}$, 17.3.1).
Avec la terminologie de la théorie de la dimension que nous introduirons au
chap. IV, il revient au même de dire qu'un \emph{préschéma localement noethérien} $T$
est de \emph{dimension cohomologique} $\leqslant n$, ou de dire qu'il est \emph{régulier}
($\mathbf{0}_{\mathrm I}$, 4.1.4) et de \emph{dimension} $\leqslant n$; cela signifie que pour
tout ouvert affine $U$ de $T$, l'anneau $\Gamma(U,\mathscr O_T)$ est \emph{de
dimension cohomologique} $\leqslant n$ ($\mathbf{0}_{\mathrm{IV}}$, 17.2.6). Cela étant,
cette dernière remarque, jointe à (6.3.2), prouve que si $S$ est \emph{localement
noethérien et de dimension cohomologique finie}, les suites spectrales
$'\mathscr E(\mathscr P_\bullet,\mathscr Q_\bullet)$ et
$''\mathscr E(\mathscr P_\bullet,\mathscr Q_\bullet)$ sont \emph{birégulières}.

Il est clair que
$\boldsymbol{\mathscr{T}\!or}_\bullet^S(\mathscr P_\bullet,\mathscr Q_\bullet)$ se
transforme en
$\boldsymbol{\mathscr{T}\!or}_\bullet^S(\mathscr Q_\bullet,\mathscr P_\bullet)$ (à un
isomorphisme près) par l'isomorphisme canonique de $X\times_S Y$ sur
$Y\times_S X$.
\end{env}

\begin{proposition}[6.5.6]
Soit $(\mathscr P_{\alpha\bullet})$ un système inductif filtrant de complexes
de $\mathscr O_X$-Modules quasi-cohérents; il existe alors un isomorphisme
canonique
\[
  \varinjlim
  (\boldsymbol{\mathscr{T}\!or}_\bullet^S
  (\mathscr P_{\alpha\bullet},\mathscr Q_\bullet))
  \xrightarrow{\sim}
  \boldsymbol{\mathscr{T}\!or}_\bullet^S
  \left(\varinjlim\mathscr P_{\alpha\bullet},\mathscr Q_\bullet\right).
  \tag{6.5.6.1}\label{III.6.5.6.1.fr}
\]
\end{proposition}

\begin{proof}
La question étant locale sur $S$, $X$ et $Y$, on peut supposer $S$, $X$, $Y$
affines et la proposition se réduit alors à (6.3.6).
\end{proof}

\begin{namedtheorem}[6.5.7]{Remarques}
\textnormal{(i)} Considérons en particulier le cas où $S=X=Y$,
$\mathscr P_\bullet$ et $\mathscr Q_\bullet$ étant donc deux complexes de
$\mathscr O_S$-Modules quasi-cohérents; alors les
$\boldsymbol{\mathscr{T}\!or}_n^S(\mathscr P_\bullet,\mathscr Q_\bullet)$ sont des
$\mathscr O_S$-Modules quasi-cohérents; en outre, pour tout point $z\in S$, il
résulte de (6.5.6) que l'on a un isomorphisme canonique
\[
  \bigl(\boldsymbol{\mathscr{T}\!or}_n^S
    (\mathscr P_\bullet,\mathscr Q_\bullet)\bigr)_z
  \xrightarrow{\sim}
  \mathbf{Tor}_n^{\mathscr O_z}
  \bigl((\mathscr P_\bullet)_z,(\mathscr Q_\bullet)_z\bigr)
  \tag{6.5.7.1}\label{III.6.5.7.1.fr}
\]
car la question est locale et on est ramené au cas des modules, en vertu de
(6.4.1.1).

\textnormal{(ii)} On peut généraliser la définition des hypertor au cas de
\emph{deux complexes de $\mathscr O_X$-Modules} $\mathscr P_\bullet$,
$\mathscr Q_\bullet$ sur un même \emph{espace annelé} $(X,\mathscr O_X)$; pour tout
ouvert $U$ de $X$, posons en effet
$A(U)=\Gamma(U,\mathscr O_X)$,
$P_\bullet(U)=\Gamma(U,\mathscr P_\bullet)$,
$Q_\bullet(U)=\Gamma(U,\mathscr Q_\bullet)$; les $A(U)$-modules
$\mathbf{Tor}_n^{A(U)}(P_\bullet(U),Q_\bullet(U))$
forment alors un \emph{préfaisceau} sur $X$, et l'on désigne par
$\boldsymbol{\mathscr{T}\!or}_n^X(\mathscr P_\bullet,\mathscr Q_\bullet)$ le
$\mathscr O_X$-Module associé à ce préfaisceau. Lorsque $X$ est un \emph{préschéma},
il résulte de (6.3.12) que ce $\mathscr O_X$-Module est canoniquement isomorphe
à l'hypertor défini ci-dessus. Nous ne développerons pas davantage cette
généralisation.
\end{namedtheorem}

\begin{proposition}[6.5.8]
Soient $X$, $Y$ deux $S$-préschémas, $\mathscr F$ \textup{(resp.} $\mathscr G$\textup{)} un
$\mathscr O_X$-Module \textup{(resp.} un $\mathscr O_Y$-Module\textup{)} quasi-cohérent. Si
$\mathscr F$ ou $\mathscr G$ est $S$-plat, on a
$\mathscr{T}\!or_n^S(\mathscr F,\mathscr G)=0$ pour $n\neq0$.
\end{proposition}

\begin{proof}
La question étant locale sur $X$ et $Y$, on peut supposer $X$, $Y$ et $S$
affines, d'anneaux respectifs $B$, $C$, $A$, et
$\mathscr F=\widetilde M$, $\mathscr G=\widetilde N$, $M$ (resp. $N$) étant
un $B$-module (resp. un $C$-module). Supposons par exemple que $\mathscr F$
soit $S$-plat, ce qui signifie que pour tout $s\in S$, $M_s$ est un
$A_s$-module plat ($\mathbf{0}_{\mathrm I}$, 6.7.1); par suite $M$ est un $A$-module plat
($\mathbf{0}_{\mathrm I}$, 6.3.3), et l'on sait que
$\operatorname{Tor}_n^A(M,N)=0$ pour $n>0$ et pour tout $C$-module $N$
($\mathbf{0}_{\mathrm I}$, 6.1.1), d'où la conclusion par (6.4.1.1).
\end{proof}

\begin{namedtheorem}[6.5.9]{Corollaire}
Soient $X$, $Y$ deux $S$-préschémas, $\mathscr P_\bullet$ \textup{(resp.}
$\mathscr Q_\bullet$\textup{)} un complexe de
\oldpage[III]{151}
$\mathscr O_X$-Modules \textup{(resp.} de $\mathscr O_Y$-Modules\textup{)} quasi-cohérents limité
inférieurement. Supposons que tous les $\mathscr P_i$ soient $S$-plats. Alors
il existe un isomorphisme canonique de $\partial$-foncteurs en
$\mathscr Q_\bullet$
\[
  \boldsymbol{\mathscr{T}\!or}_\bullet^S(\mathscr P_\bullet,\mathscr Q_\bullet)
  \xrightarrow{\sim}
  \mathscr H_\bullet(\mathscr P_\bullet\otimes_S\mathscr Q_\bullet).
  \tag{6.5.9.1}\label{III.6.5.9.1.fr}
\]
\end{namedtheorem}

\begin{proof}
Ce n'est autre que (6.3.7) lorsque $S$, $X$, $Y$ sont affines; on passe de là
au cas général par les raisonnements de (6.5.2) et (6.5.3).
\end{proof}

\begin{namedtheorem}[6.5.10]{Corollaire}
Supposons que $X$ soit plat sur $S$ \textup{($\mathbf{0}_{\mathrm I}$, 6.7.1)}, que
$\mathscr P_\bullet$ et $\mathscr Q_\bullet$ soient limités inférieurement et
que tous les $\mathscr P_i$ soient des $\mathscr O_X$-Modules localement
libres (non nécessairement de type fini). Alors l'homomorphisme \textup{(6.5.9.1)} est
bijectif.
\end{namedtheorem}

\begin{proof}
En effet, l'hypothèse de (6.5.9) est remplie, la platitude étant une propriété
ponctuelle sur $X$ par définition et toute somme directe de modules plats étant
un module plat ($\mathbf{0}_{\mathrm I}$, 6.1.2).
\end{proof}

\begin{proposition}[6.5.11]
Soient $X'$, $Y'$ deux $S$-préschémas, $f:X\to X'$, $g:Y\to Y'$ deux
$S$-morphismes affines. Soit $\mathscr P_\bullet$ \textup{(resp.}
$\mathscr Q_\bullet$\textup{)} un complexe de $\mathscr O_X$-Modules \textup{(resp.} de
$\mathscr O_Y$-Modules\textup{)} quasi-cohérents; on a alors un isomorphisme canonique
fonctoriel
\[
  (f\times_S g)_*
  \bigl(\boldsymbol{\mathscr{T}\!or}_\bullet^S
    (\mathscr P_\bullet,\mathscr Q_\bullet)\bigr)
  \xrightarrow{\sim}
  \boldsymbol{\mathscr{T}\!or}_\bullet^S
  (f_*(\mathscr P_\bullet),g_*(\mathscr Q_\bullet)).
  \tag{6.5.11.1}\label{III.6.5.11.1.fr}
\]
\end{proposition}

\begin{proof}
Comme $f$ et $g$ sont affines, $f_*(\mathscr P_\bullet)$ et
$g_*(\mathscr Q_\bullet)$ sont des complexes de Modules quasi-cohérents
(\textbf{II}, 1.2.6), et si l'on pose $Z'=X'\times_S Y'$, les deux membres de (6.5.11.1)
sont des $\mathscr O_{Z'}$-Modules quasi-cohérents (6.5.1); on se ramène
aisément au cas où $S$, $X'$ et $Y'$ sont affines; mais alors il en est de même
par hypothèse de $X$ et $Y$ et la vérification résulte aussitôt de (6.4.1.1) et
(\textbf{I}, 1.6.3).
\end{proof}

\begin{namedtheorem}[6.5.12]{Remarque}
Soient $X'$, $Y'$ deux $S$-préschémas et supposons que $X'$ soit $S$-\emph{plat};
soient $X$ un sous-préschéma fermé de $X'$, $i:X\to X'$ l'injection canonique,
$\mathscr P_\bullet$ (resp. $\mathscr Q_\bullet$) un complexe de
$\mathscr O_X$-Modules (resp. de $\mathscr O_{Y'}$-Modules) quasi-cohérents,
limité inférieurement. Soit enfin $\mathscr L'_{\bullet\bullet}$ une résolution
de $i_*(\mathscr P_\bullet)$ formé de $\mathscr O_{X'}$-Modules localement
libres, telle que tout point de $X'$ ait un voisinage ouvert affine $U$ pour
lequel $\mathscr L'_{j,\bullet}|U$ soit une résolution libre de
$i_*(\mathscr P_j)|U$ pour tout $j$. On a alors un isomorphisme canonique
\[
  (i\times_S 1)_*
  \bigl(\boldsymbol{\mathscr{T}\!or}_\bullet^{\mathbf S}
    (\mathscr P_\bullet,\mathscr Q_\bullet)\bigr)
  \xrightarrow{\sim}
  \mathscr H_\bullet
  (\mathscr L'_{\bullet\bullet}\otimes_S\mathscr Q_\bullet).
  \tag{6.5.12.1}\label{III.6.5.12.1.fr}
\]

Si $S$, $X'$, $Y'$ sont affines et si $\mathscr L'_{j,\bullet}$ est une
résolution libre de $i_*(\mathscr P_j)$ pour tout $j$, on est ramené, en vertu
de (6.5.11), au cas où $X'=X$ est $S$-plat, et il suffit d'appliquer (6.3.4).
Dans le cas général, on définit localement l'isomorphisme (6.5.12.1), et il
s'agit de vérifier que cette définition donne bien un isomorphisme global.
Pour cela, il faut se reporter à la définition du premier isomorphisme
(6.3.4.1) qui provient d'un isomorphisme de suites spectrales (\textbf{0}, 11.6.5 et
11.5.3), obtenu lui-même à partir d'un morphisme de bicomplexes
$L_{\bullet\bullet}\to L''_{\bullet\bullet}$, où $L_{\bullet\bullet}$ est la
résolution donnée de $P_\bullet$, $L''_{\bullet\bullet}$ une résolution
projective de $P_\bullet$ dans la catégorie des complexes de $A$-modules
limités
\oldpage[III]{152}
inférieurement (cf. (0, 11.5.2.2)); notre assertion résulte de ce que
l'isomorphisme (6.3.4.1) ne dépend pas de la résolution projective
$L''_{\bullet\bullet}$ choisie, en vertu de l'existence d'un homotopisme entre
deux telles résolutions (M, V, 1.2).
\end{namedtheorem}

\begin{proposition}[6.5.13]
Soient $X$, $Y$ deux $S$-préschémas, et supposons vérifiée l'une des conditions
suivantes :
\begin{enumerate}
  \item[(i)] $X$ et $Z=X\times_S Y$ sont localement noethériens et $X$ est plat
    sur $S$.
  \item[(ii)] $S$ et $X$ sont localement noethériens et $Y$ est de type fini
    sur $S$.
\end{enumerate}
Soit $\mathscr P_\bullet$ (resp. $\mathscr Q_\bullet$) un complexe de
$\mathscr O_X$-Modules (resp. de $\mathscr O_Y$-Modules) quasi-cohérents,
limité inférieurement. On suppose en outre que, pour tout $n$,
$\mathscr H_n(\mathscr P_\bullet)$ (resp.
$\mathscr H_n(\mathscr Q_\bullet)$) est un $\mathscr O_X$-Module (resp. un
$\mathscr O_Y$-Module) de type fini. Alors les
$\mathscr{T}\!or_n^S(\mathscr P_\bullet,\mathscr Q_\bullet)$ sont des
$\mathscr O_Z$-Modules cohérents.
\end{proposition}

\begin{proof}
Comme $\mathscr P_\bullet$ et $\mathscr Q_\bullet$ sont limités
inférieurement, la suite spectrale
$''\mathscr E(\mathscr P_\bullet,\mathscr Q_\bullet)$ est birégulière (6.5.4),
et en vertu de (0, 11.1.8), il suffit (puisque dans les deux cas (i), (ii), $Z$
est localement noethérien) de prouver que les termes
$''\mathscr E_{pq}^2$ sont cohérents. L'hypothèse sur les
$\mathscr H_n(\mathscr P_\bullet)$ et
$\mathscr H_n(\mathscr Q_\bullet)$ et l'expression (6.5.4.2) des
$''\mathscr E_{pq}^2$ montrent donc que la proposition est équivalente à son
cas particulier correspondant à $\mathscr P_\bullet$ et
$\mathscr Q_\bullet$ réduits à leurs termes de degré 0, autrement dit à son
\end{proof}

\begin{corollary}[6.5.14]
Supposons vérifiée l'une des conditions (i), (ii) de (6.5.13), et soit
$\mathscr F$ (resp. $\mathscr G$) un $\mathscr O_X$-Module (resp. un
$\mathscr O_Y$-Module) quasi-cohérent de type fini; alors les
$\mathscr{T}\!or_n^S(\mathscr F,\mathscr G)$ sont des $\mathscr O_Z$-Modules
cohérents.
\end{corollary}

\begin{proof}
La question étant locale sur $X$ et $Y$, on peut supposer $S$, $X$, $Y$
affines.

\textnormal{(i)} Sous les hypothèses de (i), $S$, $X$ et $Z$ sont noethériens.
Il existe donc une résolution localement libre $\mathscr L_\bullet$ de
$\mathscr F$ formée de $\mathscr O_X$-Modules de type fini (I, 1.3.7); comme
$X$ est plat sur $S$, il résulte de (6.3.4) que l'on a
\[
  \mathscr{T}\!or_n^S(\mathscr F,\mathscr G)
  =\mathscr H_n(\mathscr L_\bullet\otimes_S\mathscr G);
\]
or, les $\mathscr L_i\otimes_S\mathscr G$ sont des $\mathscr O_Z$-Modules
quasi-cohérents de type fini (I, 9.1.1), donc cohérents. On en conclut que
$\mathscr H_n(\mathscr L_\bullet\otimes_S\mathscr G)$ est cohérent
($0_{\mathrm I}$, 5.3.4).

\textnormal{(ii)} Supposons maintenant vérifiées les conditions (ii). Comme
l'anneau $A(Y)$ est quotient d'une $A(S)$-algèbre de polynômes à un nombre fini
d'indéterminées (I, 6.3.3), $Y$ est un sous-$S$-préschéma fermé d'un
$S$-préschéma affine $Y'$, plat et de type fini sur $S$; $Y'$ étant noethérien
(I, 6.3.7), il existe une résolution localement libre $\mathscr M_\bullet$ de
$j_*(\mathscr G)$ par des $\mathscr O_{Y'}$-Modules de type fini
($j:Y\to Y'$ étant l'injection canonique); en vertu de (6.5.12),
$\mathscr{T}\!or_n^S(\mathscr F,\mathscr G)$ est l'image réciproque par
$1\times j$ du $\mathscr O_{Z'}$-Module
$\mathscr H_n(\mathscr F\otimes_S\mathscr M_\bullet)$ (où
$Z'=X\times_S Y'$); on voit comme dans (i) que les
$\mathscr F\otimes_S\mathscr M_i$ sont des $\mathscr O_{Z'}$-Modules
cohérents, et l'on en tire encore la conclusion par ($0_{\mathrm I}$, 5.3.4).
\end{proof}

\begin{env}[6.5.15]
La théorie développée ci-dessus pour deux complexes $\mathscr P_\bullet$,
$\mathscr Q_\bullet$ de faisceaux quasi-cohérents sur deux $S$-préschémas
$X$, $Y$ se généralise sans peine au cas où l'on considère un nombre fini
quelconque de $S$-préschémas $X^{(i)}$ ($1\leqslant i\leqslant m$) et sur
chaque $X^{(i)}$ un complexe $\mathscr P_\bullet^{(i)}$ de
$\mathscr O_{X^{(i)}}$-Modules quasi-cohérents; si
$Z=X^{(1)}\times_S X^{(2)}\times_S\cdots\times_S X^{(m)}$, on définit ainsi un
$\mathscr O_Z$-Module quasi-cohérent
$\mathscr{T}\!or_q^S(\mathscr P_\bullet^{(1)},\mathscr P_\bullet^{(2)},
\ldots,\mathscr P_\bullet^{(m)})$. Nous laissons au lecteur le soin de
développer la théorie dans ce cas général et nous nous bornerons à écrire,
pour référence ultérieure, le terme $E_2$ de la seconde suite spectrale
(régulière) dont l'aboutissement est
$\mathscr{T}\!or_\bullet^S(\mathscr P_\bullet^{(1)},
\mathscr P_\bullet^{(2)},\ldots,\mathscr P_\bullet^{(m)})$ :

\oldpage[III]{153}

\[
  ''\mathscr E_{pq}^2
  =\bigoplus_{q_1+q_2+\cdots+q_m=q}
   \mathscr{T}\!or_p^S
   \bigl(\mathscr H_{q_1}(\mathscr P_\bullet^{(1)}),\ldots,
         \mathscr H_{q_m}(\mathscr P_\bullet^{(m)})\bigr).
  \tag{6.5.15.1}\label{III.6.5.15.1.fr}
\]
Nous étudierons dans (6.8) les suites spectrales d'associativité auxquelles
donnent lieu ces foncteurs hypertor d'un nombre quelconque de complexes.
\end{env}

